    \section{REPRESENTACIÓN GEOMÉTRICA DE LAS PARTICIONES}

    Existe un método simple para representar particiones geométricamente utilizando una
    disposición de puntos reticulares que llamaremos \textit{grafo}. Por ejemplo, la partición
    de 15 dada por
    \begin{equation*}
        6+3+3+2+1
    \end{equation*}
    se puede representar por medio de 15 puntos colocados en cinco filas tal como sigue:
    \begin{equation*}
        \begin{array}{cccccc}
            \bullet&\bullet&\bullet&\bullet&\bullet&\bullet\\
            \bullet&\bullet&\bullet&&&\\
            \bullet&\bullet&\bullet&&&\\
            \bullet&\bullet&&&&\\
            \bullet&&&&&
        \end{array}
    \end{equation*}
    Si leemos este grafo verticalmente obtenemos otra partición de 15,
    \begin{equation*}
        5+4+3+1+1+1.
    \end{equation*}
    Estas dos particiones son \textit{conjugadas}. Nótese que la parte \textit{más larga} de una
    de ellas es igual al \textit{número de partes} de la otra. Entonces tenemos el teorema que
    sigue.

    \begin{teorema}\label{teo:14-1}
        El número de particiones de $n$ en $m$ partes es igual al número de particiones de $n$
        en partes tales que la mayor sea $m$.
    \end{teorema}

    Se pueden demostrar muchos teoremas combinando simplemente argumentos que involucren
    grafos, y más adelante daremos una ilustración muy bonita de este método. Sin embargo, los
    resultados profundos de la teoría de las particiones requieren un tratamiento analítico más
    profundo, al que ahora vamos a dedicarnos.

    \section{FUNCIONES GENERADORAS DE PARTICIONES}

    Una función $F(s)$ definida por una serie de Dirichlet $F(s)=\sum f(n)n^{-s}$ se llama
    \textit{función generadora} de los coeficientes $f(n)$. Las series de Dirichlet son útiles
    para generar funciones en la teoría multiplicativa de números en virtud de la relación
    \begin{equation*}
        n^{-s}m^{-s}=(nm)^{-s}.
    \end{equation*}
    En la teoría aditiva de números es más conveniente generar funciones utilizando su
    representación en serie de potencias,
    \begin{equation*}
        F(x)=\sum f(n)x^n
    \end{equation*}
    ya que $x^nx^m=x^{n+m}$. El teorema que sigue pone de manifiesto una función generadora de
    la función de partición $p(n)$.

    \begin{teorema}[Euler]\label{teo:14-2}
        Para $\abs{x}<1$ tenemos
        \begin{equation*}
            \prod_{m=1}^{\infty}\frac{1}{1-x^m}=\sum_{n=0}^{\infty}p(n)x^n,
        \end{equation*}
        en donde $p(0)=1$.
    \end{teorema}

    \begin{proof}
        En primer lugar damos una deducción formal de esta identidad, ignorando las cuestiones
        de convergencia, y después damos una demostración más rigurosa.

        Si cada factor del producto lo desarrollamos en serie de potencias (una serie
        geométrica) obtenemos
        \begin{equation*}
            \prod_{n=1}^{\infty}\frac{1}{1-x^n}=(1+x+x^2+\cdots)(1+x^2+x^4+\cdots)(1+x^3+x^6+\cdots)\cdots
        \end{equation*}
        Ahora multiplicamos las series de la derecha, tratándolas como si fuesen polinomios y
        agrupándolas según las potencias de $x$ a fin de obtener una serie de potencias de la
        forma
        \begin{equation*}
            1+\sum_{k=1}^{\infty}a(k)x^k.
        \end{equation*}
        Queremos demostrar que $a(k)=p(k)$. Suponemos que hemos tomado el término $x^{k_1}$ de
        la primera serie, el término $x^{2k_2}$ de la segunda, el término $x^{3k_3}$ de la
        tercera, \dots, y el término $x^{mk_m}$ de la $m$-ésima, en donde cada $k_i\geq0$. Su
        producto es
        \begin{equation*}
            x^{k_1}x^{2k_2}x^{3k_3}\cdots x^{mk_m}=x^k,
        \end{equation*}
        en donde
        \begin{equation*}
            k=k_1+2k_2+3k_3+\cdots+mk_m.
        \end{equation*}
        Este resultado se puede escribir también como sigue:
        \begin{equation*}
            k=(1+1+\cdots+1)+(2+2+\cdots+2)+\cdots+(m+m+\cdots+m),
        \end{equation*}
        en donde el primer paréntesis contiene $k_1$ unos, el segundo $k_2$ doses, y así
        sucesivamente. Esta es una partición de $k$ en sumandos positivos. Entonces, cada
        partición de $k$ producirá uno de los términos $x^k$ y, recíprocamente, cada término
        $x^k$ procede de una partición conveniente de $k$. Por consiguiente $a(k)$, el
        coeficiente de $x^k$, es igual a $p(k)$, el número de particiones de $k$.

        El razonamiento anterior no constituye una demostración rigurosa ya que se han ignorado
        las cuestiones de convergencia y además hemos multiplicado conjuntamente una infinidad
        de series geométricas, tratándolas como si fuesen polinomios. Sin embargo no es difícil
        transformar las ideas anteriores en una demostración rigurosa.

        A este fin imponemos que $x$ pertenezca al intervalo $0\leq x<1$ e introducimos dos
        funciones,
        \begin{equation*}
            F_m(x)=\prod_{k=1}^{m}\frac{1}{1-x^k},\qquad\text{y}\qquad F(x)=\prod_{k=1}^{\infty}\frac{1}{1-x^k}=\lim_{m\to\infty}F_m(x).
        \end{equation*}
        El producto que define $F(x)$ converge absolutamente si $0\leq x<1$, puesto que su
        inverso $\prod(1-x^k)$ converge absolutamente (ya que la serie $\sum x^k$ converge
        absolutamente). Obsérvese además que, para cada $x$ fijo, la sucesión $\{F_m(x)\}$ es
        creciente puesto que
        \begin{equation*}
            F_{m+1}(x)=\frac{1}{1-x^{m+1}}F_m(x)\geq F_m(x).
        \end{equation*}
        Luego $F_m(x)\leq F(x)$ para cada $x$ fijo, $0\leq x<1$, y cada $m$. Pero $F_m(x)$ es
        el producto de un número finito de series absolutamente convergentes. Por consiguiente
        también ella será una serie absolutamente convergente que podemos escribir en la forma
        \begin{equation*}
            F_m(x)=1+\sum_{k=1}^{\infty}p_m(k)x^k.
        \end{equation*}
        Aquí $p_m(k)$ es el número de soluciones de la ecuación
        \begin{equation*}
            k=k_1+2k_2+\cdots+mk_m.
        \end{equation*}
        Con otras palabras, $p_m(k)$ es el número de particiones de $k$ en partes que no exceden
        $m$. Si $m\geq k$, entonces $p_m(k)=p(k)$. Por consiguiente, en todos los casos,
        tenemos
        \begin{equation*}
            p_m(k)\leq p(k)
        \end{equation*}
        con igualdad cuando $m\geq k$. Con otras palabras, tenemos
        \begin{equation*}
            \lim_{m\to\infty}p_m(k)=p(k).
        \end{equation*}
        Ahora descomponemos las series relativas a $F_m(x)$ en dos partes,
        \begin{align*}
            F_m(x)&=\sum_{k=0}^{m}p_m(k)x^k+\sum_{k=m+1}^{\infty}p_m(k)x^k\\
            &=\sum_{k=0}^{m}p(k)x^k+\sum_{k=m+1}^{\infty}p_m(k)x^k.
        \end{align*}
        Puesto que $x\geq0$ tenemos
        \begin{equation*}
            \sum_{k=0}^{m}p(k)x^k\leq F_m(x)\leq F(x).
        \end{equation*}
        Ello demuestra que la serie $\sum_{k=0}^{\infty}p(k)x^k$ converge. Sin embargo, puesto
        que $p_m(k)\leq p(k)$, tenemos
        \begin{equation*}
            \sum_{k=0}^{\infty}p_m(k)x^k\leq\sum_{k=0}^{\infty}p(k)x^k\leq F(x)
        \end{equation*}
        luego, para cada $x$ fijo, la serie $\sum p_m(k)x^k$ converge uniformemente en $m$. Si
        hacemos que $m\to\infty$ obtenemos
        \begin{equation*}
            F(x)=\lim_{m\to\infty}F_m(x)=\lim_{m\to\infty}\sum_{k=0}^{\infty}p_m(k)x^k=\sum_{k=0}^{\infty}\lim_{m\to\infty}p_m(k)x^k=\sum_{k=0}^{\infty}p(k)x^k,
        \end{equation*}
        que demuestra la identidad de Euler para $0\leq x<1$. Por prolongación analítica la
        extendemos al disco unidad $\abs{x}<1$.
    \end{proof}

    \begin{table}[H]
        \centering
        \caption{Funciones generadoras}
        \begin{tabular}{@{}lp{5.5cm}@{}}
            \hline
            Función generadora&El número de particiones de $n$ en partes que son\\\hline
            $\displaystyle\prod_{m=1}^{\infty}\frac{1}{1-x^{2m-1}}$&impares\\[2.5ex]
            $\displaystyle\prod_{m=1}^{\infty}\frac{1}{1-x^{2m}}$&pares\\[2.5ex]
            $\displaystyle\prod_{m=1}^{\infty}\frac{1}{1-x^{m^2}}$&cuadrados\\[2.5ex]
            $\displaystyle\prod_p\frac{1}{1-x^p}$&primos\\[2.5ex]
            $\displaystyle\prod_{m=1}^{\infty}(1+x^m)$&desiguales\\[2.5ex]
            $\displaystyle\prod_{m=1}^{\infty}(1+x^{2m-1})$&impares y desiguales\\[2.5ex]
            $\displaystyle\prod_{m=1}^{\infty}(1+x^{2m})$&pares y desiguales\\[2.5ex]
            $\displaystyle\prod_{m=1}^{\infty}(1+x^{m^2})$&cuadrados diferentes\\[2.5ex]
            $\displaystyle\prod_p(1+x^p)$&primos diferentes\\\hline
        \end{tabular}
    \end{table}

    Mediante razonamientos análogos podemos obtener funciones generadoras de muchas otras
    funciones de partición. En la tabla 14.1 mencionamos sólo algunos ejemplos.

    \section{EL TEOREMA DE EULER DE LOS NÚMEROS PENTAGONALES}

    A continuación consideramos la función de partición generada por el producto
    $\prod(1-x^m)$, inversa de la función generadora de $p(n)$. Escribimos
    \begin{equation*}
        \prod_{m=1}^{\infty}(1-x^m)=1+\sum_{n=1}^{\infty}a(n)x^n.
    \end{equation*}
    Para expresar $a(n)$ como una función de partición observamos que cada partición de $n$ en
    partes desiguales produce un término $x^n$ en el segundo miembro con coeficiente $+1$ ó
    $-1$. El coeficiente es $+1$ si $x^n$ es el producto de un número par de términos, y $-1$
    en los demás casos. Por consiguiente,
    \begin{equation*}
        a(n)=p_e(n)-p_o(n),
    \end{equation*}
    en donde $p_e(n)$ es el número de particiones de $n$ en un número \textit{par} de partes
    desiguales, y $p_o(n)$ es el número de particiones en un número \textit{impar} de partes
    desiguales. Euler demostró que $p_e(n)=p_o(n)$ para todo $n$ excepto para los que
    pertenecen a un conjunto especial llamado de los \textit{números pentagonales}.

    Los números pentagonales $1,5,12,22,\dots$ ya fueron mencionados en la introducción
    histórica. Están relacionados con los pentágonos dibujados en la figura 14.1.

    \begin{figure}[H]
        \centering
        \includegraphics{figuras/figura14-1.pdf}
        \caption{}
    \end{figure}

    Estos números son también las sumas parciales de los términos de la progresión aritmética
    \begin{equation*}
        1,4,7,10,13,\dots,3n+1,\dots
    \end{equation*}
    Si $\omega(n)$ designa la suma de los $n$ primeros términos de esta progresión entonces
    \begin{equation*}
        \omega(n)=\sum_{k=0}^{n-1}(3k+1)=\frac{3n(n-1)}{2}+n=\frac{3n^2-n}{2}.
    \end{equation*}
    Los números $\omega(n)$ y $\omega(-n)=(3n^2+n)/2$ se llaman \textit{números pentagonales}.

    \begin{teorema}[Teorema de los números pentagonales de Euler]\label{teo:14-3}
        Si $\abs{x}<1$ tenemos
        \begin{align*}
            \prod_{m=1}^{\infty}(1-x^m)&=1-x-x^2+x^5+x^7-x^{12}-x^{15}+\cdots\\
            &=1+\sum_{n=1}^{\infty}(-1)^n\left\{x^{\omega(n)}+x^{\omega(-n)}\right\}=\sum_{n=-\infty}^{\infty}(-1)^nx^{\omega(n)}.
        \end{align*}
    \end{teorema}

    \begin{proof}
        Primero probaremos el resultado para $0<x<1$ y después lo extenderemos al disco
        $\abs{x}<1$ por prolongación analítica. Definimos $P_0=S_0=1$ y, para $n\geq1$,
        consideramos
        \begin{equation*}
            P_n=\prod_{r=1}^{n}(1-x^r)\qquad\text{y}\qquad S_n=1+\sum_{r=1}^{n}(-1)^r\left\{x^{\omega(r)}+x^{\omega(-r)}\right\}.
        \end{equation*}
        El producto infinito $\prod(1-x^m)$ converge, luego $P_n\to\prod(1-x^m)$ cuando
        $n\to\infty$. Probaremos (utilizando el método de Shanks \cite{shanks1951}) que
        \begin{equation}\label{eq:14-6}
            \abs{S_n-P_n}\leq nx^{n+1}.
        \end{equation}
        Puesto que $nx^{n+1}\to0$ cuando $n\to\infty$ esto prueba la identidad de Euler para
        $0\leq x<1$.

        Para probar \eqref{eq:14-6} consideramos $g(r)=r(r+1)/2$ e introducimos las sumas
        \begin{equation*}
            F_n=\sum_{r=0}^{n}(-1)^r\frac{P_n}{P_r}\,x^{rn+g(r)}.
        \end{equation*}
        En primer lugar probaremos que $F_n$ es una forma disfrazada de $S_n$. Es fácil
        comprobar que $F_1=S_1=1-x-x^2$. Por consiguiente, si vemos que
        \begin{equation*}
            F_n-F_{n-1}=S_n-S_{n-1},\qquad\text{o que}\qquad F_n-S_n=F_{n-1}-S_{n-1},
        \end{equation*}
        habremos demostrado que, para todo $n\geq1$, $F_n=S_n$. Ahora bien
        \begin{equation*}
            F_n-F_{n-1}=\sum_{r=0}^{n}(-1)^r\frac{P_n}{P_r}\,x^{rn+g(r)}-\sum_{r=0}^{n-1}(-1)^r\frac{P_{n-1}}{P_r}\,x^{r(n-1)+g(r)}.
        \end{equation*}
        En la primera suma escribimos $P_n=(1-x^n)P_{n-1}$ y separamos el término con $r=n$.
        Entonces distribuimos la diferencia $1-x^n$ y obtenemos
        \begin{align*}
            F_n-F_{n-1}&=(-1)^nx^{n^2+g(n)}+\sum_{r=0}^{n-1}(-1)^r\frac{P_{n-1}}{P_r}\,x^{rn+g(r)}\\
            &\quad-\sum_{r=0}^{n-1}(-1)^r\frac{P_{n-1}}{P_r}\,x^{(r+1)n+g(r)}-\sum_{r=0}^{n-1}(-1)^r\frac{P_{n-1}}{P_r}\,x^{r(n-1)+g(r)}.
        \end{align*}
        Ahora combinamos la primera y la tercera suma y observamos que el término con $r=0$ se
        anula. En el segundo sumando desplazamos los índices y obtenemos
        \begin{align*}
            F_n-F_{n-1}&=(-1)^nx^{n^2+g(n)}+\sum_{r=1}^{n-1}(-1)^r\frac{P_{n-1}}{P_r}\,x^{r(n-1)+g(r)}(x^r-1)\\
            &\quad-\sum_{r=1}^{n}(-1)^{r-1}\frac{P_{n-1}}{P_{r-1}}\,x^{rn+g(r-1)}.
        \end{align*}
        Pero $(x^r-1)/P_r=-1/P_{r-1}$ y $r(n-1)+g(r)=rn+g(r-1)$ luego los dos últimos sumandos
        se anulan término a término excepto para el término con $r=n$ de la segunda suma. Así,
        obtenemos
        \begin{equation*}
            F_n-F_{n-1}=(-1)^nx^{n^2+g(n)}+(-1)^nx^{n^2+g(n-1)}.
        \end{equation*}
        Pero
        \begin{equation*}
            n^2+g(n)=n^2+\frac{n(n+1)}{2}=\omega(-n)\qquad\text{y}\qquad n^2+g(n-1)=\omega(n),
        \end{equation*}
        luego
        \begin{equation*}
            F_n-F_{n-1}=(-1)^n\left\{x^{\omega(n)}+x^{\omega(-n)}\right\}=S_n-S_{n-1},
        \end{equation*}
        y entonces $F_n=S_n$ para todo $n\geq1$. En la suma que define $F_n$ el primer término
        es $P_n$ luego
        \begin{equation}\label{eq:14-7}
            F_n=P_n+\sum_{r=1}^{n}(-1)^r\frac{P_n}{P_r}\,x^{rn+g(r)}.
        \end{equation}
        Obsérvese que $0<P_n/P_r\leq1$ ya que $0\leq x<1$. Además, cada factor
        $x^{rn+g(r)}\leq x^{n+1}$, luego la suma de la derecha de \eqref{eq:14-7} está acotada
        superiormente por $nx^{n+1}$. Por consiguiente $\abs{F_n-P_n}\leq nx^{n+1}$ y, dado que
        $F_n=S_n$, esto demuestra \eqref{eq:14-6} y termina la demostración de la identidad de
        Euler.
    \end{proof}

    \section{DEMOSTRACIÓN COMBINATORIA DEL TEOREMA DE LOS NÚMEROS PENTAGONALES DE EULER}

    Euler demostró su teorema de los números pentagonales por inducción en 1750. Más tarde
    Legendre en 1830 y Jacobi en 1846 obtuvieron otras demostraciones. Esta sección describe
    una demostración combinatoria realmente notable dada por F. Franklin \cite{franklin1881} en
    1881.

    Ya hemos observado que
    \begin{equation*}
        \prod_{m=1}^{\infty}(1-x^m)=1+\sum_{n=1}^{\infty}\left\{p_e(n)-p_o(n)\right\}x^n,
    \end{equation*}
    en donde $p_e(n)$ es el número de particiones de $n$ en un número \textit{par} de partes
    desiguales, y $p_o(n)$ es el número de particiones de $n$ en un número \textit{impar} de
    partes desiguales. Franklin utilizó la representación gráfica de las particiones mediante
    puntos reticulares para demostrar que existe una correspondencia uno a uno entre las
    particiones de $n$ en un número par de partes desiguales y las particiones de $n$ en un
    número impar de partes desiguales, de modo que $p_e(n)=p_o(n)$, excepto cuando $n$ es un
    número pentagonal.

    Consideremos el grafo de una partición de $n$ partes desiguales. Decimos que el grafo se
    halla en \textit{forma estándar} si las partes se colocan en orden decreciente, tal como
    ilustra el ejemplo de la figura 14.2. La longitud del segmento rectilíneo que contiene los
    puntos de la última fila se llama la \textit{base} del grafo, y el número de puntos
    reticulares de la base se designará por $b$. Luego, $b\geq1$. El segmento rectilíneo más
    largo de $45^{\circ}$ de inclinación que une el último punto de la primera fila con otros
    puntos del grafo se denomina \textit{rampa} del grafo, y el número de puntos reticulares que
    se hallan sobre la rampa se designa por $s$. Luego $s\geq1$. En la figura 14.2 tenemos
    $b=2$ y $s=4$.

    \begin{figure}[H]
        \centering
        \includegraphics{figuras/figura14-2.pdf}
        \caption{}
    \end{figure}

    Ahora definimos dos operaciones $A$ y $B$ en este grafo. La operación $A$ mueve los puntos
    de la base de forma que pertenezcan a una recta paralela a la rampa, tal como se indica en
    la figura 14.3(a). La operación $B$ mueve los puntos de la rampa de forma que pertenezcan a
    una recta paralela a la base, tal como se indica en la figura 14.3(b). Decimos que una
    operación es \textit{permisible} si preserva la forma estándar del grafo, es decir, si el
    nuevo grafo tiene también las partes desiguales colocadas en orden decreciente.

    \begin{figure}[H]
        \centering
        \includegraphics{figuras/figura14-3.pdf}
        \caption{}
    \end{figure}

    Si $A$ es permisible obtenemos una nueva partición de $n$ en partes desiguales, pero el
    número de partes es una unidad menor que antes de efectuar la operación $A$. Si $B$ es
    permisible obtenemos una nueva partición en partes desiguales, pero el número de partes de
    $n$ es una unidad mayor que antes. Por consiguiente, si para cada partición de $n$
    exactamente una de las operaciones $A$ o $B$ es permisible, existe una correspondencia uno a
    uno entre particiones de $n$ en partes desiguales en número par e impar, es decir
    $p_e(n)=p_o(n)$ para dicho $n$.

    Para determinar cuándo $A$ y $B$ son permisibles consideramos tres casos: (1) $b<s$;
    (2) $b=s$; (3) $b>s$.

    \textit{Caso 1:} Si $b<s$, entonces $b\leq s-1$, luego la operación $A$ es permisible, pero
    la $B$ no porque destruye la forma estándar. (Ver figura 14.3.)

    \textit{Caso 2:} Si $b=s$, la operación $B$ no es permisible ya que el grafo que resulta no
    se halla en forma estándar. La operación $A$ es permisible excepto cuando la base y la
    rampa se intersecan como muestra la figura 14.4(a), en cuyo caso el nuevo grafo no se halla
    en forma estándar.

    \textit{Caso 3:} Si $b>s$, la operación $A$ no es permisible, mientras que la $B$ es
    permisible excepto si $b=s+1$ y la base y la rampa se intersecan, como muestra la figura
    14.4(b). En este caso el nuevo grafo tiene dos partes iguales.

    \begin{figure}[H]
        \centering
        \includegraphics{figuras/figura14-4.pdf}
        \caption{Ni $A$ ni $B$ son permisibles}
    \end{figure}

    Por consiguiente, exactamente una de las operaciones $A$ o $B$ es permisible, con las dos
    excepciones arriba indicadas. Consideremos en primer lugar el caso excepcional indicado en
    la figura 14.4(a) y supongamos que el grafo posee $k$ filas. Entonces también $b=k$, luego
    el número $n$ se obtiene por
    \begin{equation*}
        n=k+(k+1)+\cdots+(2k-1)=\frac{3k^2-k}{2}=\omega(k).
    \end{equation*}
    Para esta partición de $n$ tenemos una partición extra en partes pares si $k$ es par, y una
    partición extra en partes impares si $k$ es impar, es decir
    \begin{equation*}
        p_e(n)-p_o(n)=(-1)^k.
    \end{equation*}
    En el otro caso excepcional, indicado en la figura 14.4(b), existe un punto reticular
    adicional en cada fila o sea
    \begin{equation*}
        n=\frac{3k^2-k}{2}+k=\frac{3k^2+k}{2}=\omega(-k)
    \end{equation*}
    y de nuevo $p_e(n)-p_o(n)=(-1)^k$. Esto termina la demostración de Franklin de la identidad
    de Euler.

    \section[FÓRMULA RECURSIVA DE EULER PARA $p(n)$]{FÓRMULA RECURSIVA DE EULER PARA \texorpdfstring{$p(n)$}{p(n)}}

    \begin{teorema}\label{teo:14-4}
        Sea $p(0)=1$ y definamos $p(n)$ igual a 0 si $n<0$. Entonces, para $n\geq1$, tenemos
        \begin{equation}\label{eq:14-8}
            p(n)-p(n-1)-p(n-2)+p(n-5)+p(n-7)+\cdots=0,
        \end{equation}
        o bien, lo que es lo mismo,
        \begin{equation*}
            p(n)=\sum_{k=1}^{\infty}(-1)^{k+1}\left\{p(n-\omega(k))+p(n-\omega(-k))\right\}.
        \end{equation*}
    \end{teorema}

    \begin{proof}
        Los teoremas \ref{teo:14-2} y \ref{teo:14-3} nos dan la identidad
        \begin{equation*}
            \left(1+\sum_{k=1}^{\infty}(-1)^k\left\{x^{\omega(k)}+x^{\omega(-k)}\right\}\right)\left(\sum_{m=0}^{\infty}p(m)x^m\right)=1.
        \end{equation*}
        Si $n\geq1$, el coeficiente de $x^n$ de la derecha es 0, luego inmediatamente
        obtenemos \eqref{eq:14-8} igualando coeficientes.
    \end{proof}

    MacMahon utilizó esta fórmula recursiva para calcular $p(n)$ hasta $n=200$. Damos algunos
    valores muestra de esta tabla:
    \begin{center}
        \begin{tabular}{@{}l@{}}
            $p(1)=1$\\
            $p(5)=7$\\
            $p(10)=42$\\
            $p(15)=176$\\
            $p(20)=627$
        \end{tabular}
    \end{center}

